\documentclass[10pt,a4paper]{article}
\usepackage[T1]{fontenc}
\usepackage[utf8]{inputenc}
\usepackage{lmodern}
\usepackage{amsmath,amssymb}
\usepackage[margin=21mm]{geometry}
\usepackage{microtype}
\usepackage[hidelinks]{hyperref}
\usepackage{enumitem}
\setlist[itemize]{leftmargin=*,nosep,topsep=3pt}
\setlength{\parindent}{0pt}
\setlength{\parskip}{5pt}
\setlength{\emergencystretch}{3em}
\newcommand{\question}[1]{\par\medskip{\normalsize\bfseries #1}\par\smallskip}
\pagestyle{plain}

\begin{document}
{\Large\bfseries Towards adaptive deep brain stimulation for freezing of gait\par}

\textbf{Paper:} Huiling Tan (2022).\quad Scientific commentary on Yin et al. (note 02).\par
\href{https://doi.org/10.1093/brain/awac172}{DOI: \nolinkurl{10.1093/brain/awac172}}\par
{\footnotesize Page references below use PDF page numbers, including front matter.}

\question{To which general question does this paper relate?}
Feedback-signal selection and the distinction between a predictive biomarker and a causal stimulation target, especially for freezing of gait.

\question{How does it answer the general questions?}
\textbf{Argument.} The commentary discusses cortical beta--gamma PAC as a freezing-related signal and the finding that DBS can improve freezing beyond its effect on PAC. It contributes interpretation, not a new patient experiment (PDF pp.~1--3).

\textbf{Usable insight.} Statistical PAC can reflect nested oscillations or non-sinusoidal waveform shape. A PAC change should therefore not automatically be interpreted as altered interaction between two independent rhythms. Biomarker-triggered treatment and stimulation designed to directly disrupt a causal pathological process require different levels of mechanistic evidence (PDF pp.~1--3).

\textbf{Model relevance.} Useful criteria for choosing feedback: symptom relevance, reliable estimation and a justified treatment response. Freezing is the clinical outcome; PAC or beta suppression is a proposed intermediate mechanism to test. PAC estimation is not a linear state measurement, and the commentary gives no delayed-state parameterization or demonstrated PAC-driven controller.

\question{What papers does it cite that could provide more information?}
\begin{itemize}
\item Yin et al. (2022), \emph{Cortical phase-amplitude coupling is key to the occurrence and treatment of freezing of gait}: underlying experiment (ref.~2, PDF p.~3).
\item Cole et al. (2017), \emph{Nonsinusoidal beta oscillations reflect cortical pathophysiology in Parkinson's disease}: PAC interpretation (ref.~3, PDF p.~3).
\end{itemize}

\end{document}
